% ============================================================================
% MIST — Results tables (WS2 / task #27). Generated from scripts/eval_b1.py.
% Requires \usepackage{booktabs} in the preamble.
%
% Checkpoint: checkpoints/csf_b1.pt  (motion input, 15 epochs, lr 1e-4)
% All methods evaluated on the SAME (kp_a, kp_b, Δt) pairs per split.
% Accin@τ = fraction of predictions within τ frames; higher is better.
% Frm.err, MAE(ms): lower is better. FPS = 29.97.
% ============================================================================

% ---------------------------------------------------------------------------
% Table 1: Main results on the Panoptic B1 benchmark (validation + held-out test)
% ---------------------------------------------------------------------------
\begin{table}[t]
\centering
\caption{Sub-frame synchronization on the Panoptic B1 benchmark. ContinuSyncFormer
is our method; Cross-Correlation and DTW are classical baselines. All methods see
identical sample pairs. Best per column in \textbf{bold}.}
\label{tab:b1_main}
\begin{tabular}{llrrrr}
\toprule
Split & Method & Frm.err $\downarrow$ & Acc$_{in}$@0.1 $\uparrow$ & Acc$_{in}$@0.25 $\uparrow$ & MAE (ms) $\downarrow$ \\
\midrule
\multirow{3}{*}{Validation}
 & ContinuSyncFormer (ours) & \textbf{0.237} & \textbf{0.469} & \textbf{0.811} & \textbf{7.90} \\
 & CC + parabolic           & 0.430 & 0.435 & 0.746 & 14.34 \\
 & DTW                      & 0.359 & 0.429 & 0.764 & 11.96 \\
\midrule
\multirow{3}{*}{Test (held-out)}
 & ContinuSyncFormer (ours) & \textbf{0.126} & 0.556 & \textbf{0.929} & \textbf{4.21} \\
 & CC + parabolic           & 0.205 & \textbf{0.628} & 0.926 & 6.85 \\
 & DTW                      & 0.156 & 0.590 & \textbf{0.929} & 5.20 \\
\bottomrule
\end{tabular}
\end{table}

% ---------------------------------------------------------------------------
% Table 2: Occlusion robustness (validation, 20% per-keypoint occlusion).
% Deterministic masks: baselines are identical across both model runs.
% ---------------------------------------------------------------------------
\begin{table}[t]
\centering
\caption{Occlusion robustness on validation with 20\% per-keypoint occlusion; all
methods see identical occlusion masks. Occlusion-aware training (reliability-biased
attention) improves the fraction of accurate predictions (Acc$_{in}$@$\tau$) at the
cost of a slightly higher mean error. Best per column in \textbf{bold}.}
\label{tab:b1_occlusion}
\begin{tabular}{lrrrr}
\toprule
Method & Frm.err $\downarrow$ & Acc$_{in}$@0.1 $\uparrow$ & Acc$_{in}$@0.25 $\uparrow$ & MAE (ms) $\downarrow$ \\
\midrule
ContinuSyncFormer (standard)     & \textbf{0.318} & 0.342 & 0.683 & \textbf{10.59} \\
ContinuSyncFormer (occ.\ aware)  & 0.384 & \textbf{0.396} & \textbf{0.710} & 12.81 \\
CC + parabolic                   & 0.518 & 0.365 & 0.675 & 17.30 \\
DTW                              & 0.386 & 0.374 & 0.686 & 12.89 \\
\bottomrule
\end{tabular}
\end{table}

% ---------------------------------------------------------------------------
% Table 3: Ablation on B1 validation. One axis flipped per row; 6 epochs each,
% best val Acc@0.1. Single run — treat differences < ~0.03 as noise.
% ---------------------------------------------------------------------------
\begin{table}[t]
\centering
\caption{Component ablation on Panoptic B1 validation (6 epochs, best
Acc$_{in}$@0.1). Each row flips one axis of the full model. The motion-feature
input is decisive; the positional-encoding variant and cross-view attention have
negligible effect at this scale.}
\label{tab:b1_ablation}
\begin{tabular}{llrr}
\toprule
Axis & Variant & Acc$_{in}$@0.1 $\uparrow$ & $\Delta$ \\
\midrule
--- & full (RoPE + cross-view + motion) & 0.458 & --- \\
\midrule
\multirow{2}{*}{Positional enc.}
 & sincos                & 0.442 & $-0.016$ \\
 & none                  & 0.484 & $+0.026$ \\
\midrule
Cross-view attn.
 & self-only (no cross)  & 0.453 & $-0.005$ \\
\midrule
\multirow{2}{*}{Input repr.}
 & pose coords (no motion)        & 0.063 & $-0.395$ \\
 & raw coords (no motion, no norm) & 0.012 & $-0.446$ \\
\bottomrule
\end{tabular}
\end{table}

% ---------------------------------------------------------------------------
% Table 4: Downstream 3D — triangulation MPJPE vs desynchronization (B3).
% Held-out test sequences; 3 HD cameras (nodes 0/10/20); 11 clips.
% ---------------------------------------------------------------------------
\begin{table}[t]
\centering
\caption{Sub-frame desynchronization corrupts multi-view triangulation, and sync
correction recovers it. MPJPE (mm) vs the true 3D on held-out Panoptic sequences;
one camera is the reference, the others are shifted by $\pm$offset. \emph{naive}
= no correction, \emph{oracle} = corrected with the true offset. At offset 0 the
naive MPJPE is exactly 0, validating the project$\rightarrow$triangulate loop.}
\label{tab:b3_triangulation}
\begin{tabular}{rrrrr}
\toprule
Offset (fr) & naive & CC & ContinuSyncFormer & oracle \\
\midrule
0.00 & \textbf{0.00} & 0.49 & 2.03 & 0.00 \\
0.25 & 0.92 & 0.65 & 1.96 & \textbf{0.19} \\
0.50 & 1.91 & 0.99 & 2.17 & \textbf{0.36} \\
1.00 & 3.88 & 0.58 & 1.56 & \textbf{0.00} \\
2.00 & 7.64 & 0.53 & 1.51 & \textbf{0.00} \\
3.00 & 11.27 & 0.53 & 1.47 & \textbf{0.00} \\
\bottomrule
\end{tabular}
\end{table}

% ---------------------------------------------------------------------------
% Table 5: Drift recovery -- the regime where the learned method wins.
% Time-varying offset t_B = alpha*t_A + beta; 8 clips; window=20.
% ---------------------------------------------------------------------------
\begin{table}[t]
\centering
\caption{Recovering a \emph{drifting} inter-camera offset (mean~$\pm$~std over 8 clips of
$|\hat\delta(t)-\delta(t)|$ in frames; lower better). CC-const cannot track drift;
Caspi--Irani is a coarse affine grid; CC-slide degrades on the short windows drift tracking
needs and its std \emph{exceeds its mean} at high drift (it catastrophically fails on some
clips). CSF-slide wins at every non-zero drift with far lower variance. Last column: mean~$\pm$~std
over 4 training seeds -- the win is not a lucky seed.}
\label{tab:b3_drift}
\begin{tabular}{rrrrrr}
\toprule
Drift (fr) & CC-const & Caspi & CC-slide & \textbf{CSF-slide} & CSF (over seeds) \\
\midrule
0.0 & 0.19\,$\pm$\,0.19 & 0.30\,$\pm$\,0.25 & 0.18\,$\pm$\,0.18 & 0.26\,$\pm$\,0.38 & 0.17\,$\pm$\,0.07 \\
1.0 & 0.32\,$\pm$\,0.11 & 0.28\,$\pm$\,0.20 & 0.20\,$\pm$\,0.17 & \textbf{0.11\,$\pm$\,0.03} & 0.10\,$\pm$\,0.01 \\
2.0 & 0.59\,$\pm$\,0.14 & 0.24\,$\pm$\,0.11 & 0.29\,$\pm$\,0.30 & \textbf{0.10\,$\pm$\,0.07} & 0.11\,$\pm$\,0.02 \\
3.0 & 0.84\,$\pm$\,0.17 & 0.27\,$\pm$\,0.14 & 0.42\,$\pm$\,0.56 & \textbf{0.13\,$\pm$\,0.11} & 0.12\,$\pm$\,0.02 \\
\bottomrule
\end{tabular}
\end{table}

% ---------------------------------------------------------------------------
% Table 6: Downstream 3D under drift -- MPJPE (mm). 2 clips (illustrative), window=20.
% ---------------------------------------------------------------------------
\begin{table}[t]
\centering
\caption{Triangulation MPJPE (mm) under a drifting offset. A constant correction
(CC-const) lets 3D error grow with drift; the learned sliding-window correction holds
the reconstruction nearest the perfect-sync oracle across the take.}
\label{tab:b3_drift_mpjpe}
\begin{tabular}{rrrrrr}
\toprule
Drift (fr) & naive & CC-const & CC-slide & \textbf{CSF-slide (ours)} & oracle \\
\midrule
0.0 & 3.65 & 1.63 & 1.69 & \textbf{1.55} & 0.37 \\
0.5 & 5.32 & 1.60 & 1.75 & \textbf{1.36} & 0.21 \\
1.0 & 6.93 & 1.90 & 2.16 & \textbf{1.49} & 0.20 \\
2.0 & 10.10 & 3.37 & 2.65 & \textbf{1.60} & 0.25 \\
3.0 & 13.20 & 4.93 & 3.52 & \textbf{1.72} & 0.36 \\
\bottomrule
\end{tabular}
\end{table}

% ---------------------------------------------------------------------------
% Table 7: Drift recovery vs window size (mechanism). T=20 model, 8 clips.
% CSF-slide is robust at every window; CC-slide is reliable only at long windows.
% ---------------------------------------------------------------------------
\begin{table}[t]
\centering
\caption{Drift-recovery error (frames) vs sliding-window length, at 2-frame drift.
CC-slide needs long windows to be reliable and still degrades; CSF-slide is robust
down to very short windows -- which is what enables fine, low-latency drift tracking.}
\label{tab:b3_window}
\begin{tabular}{rrr}
\toprule
Window (frames) & CC-slide & \textbf{CSF-slide (ours)} \\
\midrule
16 & 0.268 & \textbf{0.091} \\
20 & 0.288 & \textbf{0.101} \\
24 & 0.207 & \textbf{0.116} \\
32 & \textbf{0.101} & 0.184 \\
48 & \textbf{0.145} & 0.140 \\
\bottomrule
\end{tabular}
\end{table}

% ---------------------------------------------------------------------------
% Table 8: Non-linear (oscillating) drift. window=16, period=100, 3 seqs.
% Only the learned estimator can track a curve; classical curve-tracking is worse
% than predicting the mean because CC's short-window estimates are too noisy.
% ---------------------------------------------------------------------------
\begin{table}[t]
\centering
\caption{Recovering a \emph{non-linear} (sinusoidal) offset (mean error, frames).
A line/affine fit cannot represent a curve, so CC-const, Caspi--Irani and CC-line
just predict the mean -- their error grows with amplitude. Classical curve tracking
(CC-curve) is even worse: CC's short-window estimates are too noisy. Only the learned
curve tracker (CSF-curve) keeps the error nearly amplitude-independent.}
\label{tab:b3_nonlinear}
\begin{tabular}{rrrrrr}
\toprule
Amp (fr) & CC-const & Caspi & CC-line & CC-curve & \textbf{CSF-curve (ours)} \\
\midrule
0.5 & 1.75 & 1.37 & \textbf{0.30} & 0.89 & 0.56 \\
1.0 & 2.05 & 1.60 & \textbf{0.59} & 0.92 & 0.64 \\
2.0 & 2.70 & 2.16 & 1.18 & 1.33 & \textbf{0.73} \\
\bottomrule
\end{tabular}
\end{table}
